\begin{frame}{Методология: трассировки Acme}
	Анализ основан на данных за 6 месяцев (март–август 2023) из двух кластеров для LLM:
	\begin{itemize}
		\item \textbf{Seren:} 2288 GPU A100;
		\item \textbf{Kalos:} 2416 GPU A100.
	\end{itemize}
	
	Источники данных:
	\begin{itemize}
		\item журналы планировщика задач;
		\item мониторинг инфраструктуры (Prometheus, NVIDIA DCGM);
		\item журналы выполнения и профилирование.
	\end{itemize}
	
	Трассировки охватывают сотни тысяч заданий CPU/GPU и детальные метрики оборудования.
\end{frame}

\begin{frame}{Ключевые наблюдения — время выполнения и задержки}
	\begin{itemize}
		\item \textbf{Короткая продолжительность заданий:} 
		Средняя длительность заданий LLM в 2,7–12,8 раз короче, чем в ранее опубликованных DL‑трассировках. 
		Это связано с множеством краткосрочных задач, таких как оценка и тюнинг.
		
		\item \textbf{Несправедливая задержка в очереди:} 
		Короткие оценочные задания испытывают наибольшие задержки в очереди, 
		потому что ресурсы преимущественно резервируются для долгих предобучающих задач (pretraining).
	\end{itemize}
\end{frame}

\begin{frame}{Ключевые наблюдения — использование ресурсов}
	\textbf{Несбалансированное потребление:}
	\begin{itemize}
		\item Предобучение (pretraining) — 3.2\% заданий, но потребляет 94.0\% GPU‑времени (Kalos);
		\item Оценка (evaluation) — 92.9\% заданий, но использует только 0.8\% ресурсов.
	\end{itemize}
	
	\textbf{Недоиспользование сопутствующих ресурсов:}
	\begin{itemize}
		\item CPU, хост‑память и сеть часто не загружены;
		\item GPU и память GPU показывают высокие медианные значения:
		\begin{itemize}
			\item Kalos: GPU‑память медиана 75\% (60 ГБ), GPU‑использование 99\%
		\end{itemize}
	\end{itemize}
\end{frame}

\begin{frame}{Проблемы в оценочных заданиях}
	Профилирование показало значительное время ожидания и простоя GPU в оценочных конвейерах:
	\begin{itemize}
		\item Пример: HumanEval тратит 29.5\% времени на загрузку модели и предобработку данных;
		\item 19.0\% времени уходит на проверку корректности синтезированных программ;
		\item Лишь около 50\% времени используется непосредственно для вывода на GPU.
	\end{itemize}
	
	\textbf{Следствие:} это удлиняет очереди при массовой оценке и увеличивает задержки выполнения заданий.
\end{frame}

\begin{frame}{Системные решения Acme — обзор}
	На основе наблюдений разработаны две ключевые системы:
	\begin{enumerate}
		\item \textbf{Отказоустойчивое предварительное обучение:} для диагностики сбоев и автоматического восстановления.
		\item \textbf{Раздельное планирование оценки:} для ускорения обратной связи и сокращения простоя GPU.
	\end{enumerate}
\end{frame}

\begin{frame}{Отказоустойчивое предварительное обучение}
	\begin{itemize}
		\item Частые асинхронные контрольные точки (checkpoints) для быстрого восстановления.
		\item Диагностика корневых причин с помощью эвристик и LLM‑помощника.
		\item Инструменты для обнаружения неисправных узлов и автоматического перезапуска из корректной контрольной точки.
	\end{itemize}
	\textbf{Эффект:} ускорение проверки восстановлений в 3.6–58.7 раза и значительное сокращение ручного вмешательства.
\end{frame}

\begin{frame}{Раздельное планирование оценки и выводы}
	\textbf{Развязка извлечения модели:}
	\begin{itemize}
		\item Избегает конфликтов при удалённой загрузке модели за счёт отдельной фазы извлечения.
	\end{itemize}
	
	\textbf{Развязка вычисления метрик:}
	\begin{itemize}
		\item Минимизирует простой GPU, параллелит вычисление метрик и балансирует нагрузку по предзнаниям и наборам данных.
	\end{itemize}
	
	\textbf{Практический эффект:}
	\begin{itemize}
		\item Сокращение времени оценки до 1.8x.
		\item Повышена своевременная обратная связь по качеству модели.
	\end{itemize}
	
	\textbf{Ресурсы:} трассировки и код опубликованы на \url{https://github.com/InternLM/AcmeTrace}.
\end{frame}